\subsection{Projective motives and Lefschetz symmetry for the actual word traces}
The whole-word trace from \fid{F43} has a stronger realization than
coefficientwise positivity. An elementary symmetric-square identity
and a rational algebraic projector supply a smooth-projective
coefficient motive with a specified Lefschetz operator. The motive
depends on the word; no common operator realization is assumed.

\begin{finding}[\fid{F45}: a projective word-trace realization]\label{f:projectiveword}
Let $w$ be a nonempty word in $P_1,P_2$, of length $L$ and
cyclic transition count $2s$. Thus $0\le2s\le L$, and its actual
trace is
\[
 T_w(v)=20D^L+16D^{L-s}+4D^{L-2s}(D-2)^{2s},
 \qquad D=(v+v^{-1})^2.
\]
This is the character of an actual finite-dimensional
$\mathrm{SL}_2$-representation with even weights. Its coefficients
are symmetric and unimodal in $v^2$.

Over every $k=\F_q$, with $\ell\ne\operatorname{char}k$, there is
an explicit rational Chow-motive projector on products of
projective lines whose centered realization $\mathcal K_w$ has
\begin{equation}\label{eq:projectivewordtrace}
 \sum_i(-1)^i\Tr(F^n\mid H^i(\mathcal K_w))
             =T_w(q^{n/2})\qquad(n\ge1).
\end{equation}
Only even degrees occur, and if $T_w(v)=\sum_r c_rv^{2r}$, then
\[
 H^{2r}(\mathcal K_w)=\Q_\ell(-r)^{\oplus c_r}.
\]
Thus $\mathcal K_w$ is pure of weight zero as a complex in BBD's
convention, while its total ungraded Weil space is mixed.
An invariant ample class gives hard Lefschetz on the projector
image and on the centered direct sum.
\end{finding}
\begin{proof}
\par\smallskip\noindent\textbf{An actual character.}\quad
Let $V$ be the standard two-dimensional characteristic-zero
$\mathrm{SL}_2$-representation, with character $\delta=v+v^{-1}$,
and put $W=V^{\otimes2s}$. For eigenvalues $z_i$, the identity
\[
 2\sum_{i\le j}z_iz_j=(\sum_i z_i)^2+\sum_i z_i^2
\]
proves
\[
 \chi_{\Sym^2W}(v)
   =\tfrac12\bigl(D^{2s}+(v^2+v^{-2})^{2s}\bigr)
   =\tfrac12\bigl(D^{2s}+(D-2)^{2s}\bigr).
\]
All tensor exponents in the actual representation
\begin{align*}
 \mathcal R_w={}&(V^{\otimes2L})^{\oplus16}
 \oplus(V^{\otimes2(L-s)})^{\oplus16}\\
 &\oplus\left(V^{\otimes(2L-4s)}
                    \otimes\Sym^2(V^{\otimes2s})\right)^{\oplus8}
\end{align*}
are nonnegative. Its character is
\[
 16D^L+16D^{L-s}
 +8D^{L-2s}\tfrac12\bigl(D^{2s}+(D-2)^{2s}\bigr)=T_w.
\]
All constituents have even highest weights. Their characters are
the centered strings $v^{-2a}+v^{-2a+2}+\cdots+v^{2a}$.
A nonnegative sum of these strings is symmetric and nondecreasing
toward the center in $v^2$. This proves the first assertion.
For $s=0$, the zero-fold tensor power and its symmetric square
are the unit, and the formula gives $40D^L$.

\par\smallskip\noindent\textbf{The projective schemes and projector.}\quad
For $d\ge0$ put
\[
 X_d=(\PP^1_k)^{2d},\qquad
 K_d=R\Gamma(X_{d,\bar k},\Q_\ell)(d)[2d].
\]
For $d=0$ these are the point and unit complex. The scheme $X_d$
is smooth projective of dimension $2d$.
On $X_L$, exchange two ordered blocks of $2s$ factors and leave
the other $2L-4s$ factors fixed. This defines a $k$-rational
involution $\sigma$, with $2s$ transposed pairs; when $s=0$ it
is the identity. Its graph gives the actual rational correspondence
\[
 \pi_s=\tfrac12(\Delta_{X_L}+\Gamma_\sigma)
       \in CH^{2L}(X_L\times X_L)_\Q.
\]
Since $\Gamma_\sigma^2=\Delta_{X_L}$, this is an idempotent.
Its image exists in the idempotent-complete category of rational
Chow motives and realizes the invariant projector $(1+\sigma^*)/2$.
No smoothness of a quotient variety is used. Rational coefficients
also permit characteristic two: the scalar one-half is in the
characteristic-zero coefficient field.

Define
\begin{equation}\label{eq:wordcomplex}
 \mathcal K_w=K_L^{\oplus16}\oplus K_{L-s}^{\oplus16}
                         \oplus(\pi_sK_L)^{\oplus8}.
\end{equation}
Before the centering shifts, this comes from the rational motives
$h(X_L)(L)$, $h(X_{L-s})(L-s)$ and the image of $\pi_s$ on
$h(X_L)(L)$, with multiplicities $16,16,8$.
The first and third realizations are shifted by $2L$, the second
by $2(L-s)$, exactly as displayed. All Tate twists are integers;
the notation $q^{n/2}$ only evaluates an even Laurent polynomial.

\par\smallskip\noindent\textbf{All extension traces.}\quad
The projective-space and projective-bundle calculations
\cite[Example~16.3, p.~109; Theorem~23.2, pp.~139--140]{Milne}
give the product basis
\[
 H^*(X_{d,\bar k})=\Q_\ell[x_1,\ldots,x_{2d}]/(x_i^2),
 \qquad \deg x_i=2.
\]
The monomial $x_I=\prod_{i\in I}x_i$ spans a Tate line with
geometric Frobenius $F^n$ eigenvalue $Q^{|I|}$, where $Q=q^n$.
There are no odd groups. The even shift changes no Euler sign,
and consequently
\[
 \Tr(F^n\mid K_d)=Q^{-d}(1+Q)^{2d}=D_n^d,
 \qquad D_n=Q+2+Q^{-1}.
\]
Here trace of a complex means its finite alternating trace.

The involution permutes the subset basis without Koszul signs,
since all factor classes have even degree. To be fixed, a subset
may choose each fixed coordinate independently and must choose
both or neither coordinate in each transposed pair. Thus
\[
 \Tr(\sigma F^n\mid K_L)
   =Q^{-L}(1+Q)^{2L-4s}(1+Q^2)^{2s}
   =D_n^{L-2s}(D_n-2)^{2s}.
\]
The projector commutes with Frobenius, so its actual image has trace
\[
 \Tr(F^n\mid\pi_sK_L)
   =\tfrac12\left[D_n^L+D_n^{L-2s}(D_n-2)^{2s}\right].
\]
Combining the $16,16,8$ summands gives precisely
$20D_n^L+16D_n^{L-s}+4D_n^{L-2s}(D_n-2)^{2s}$,
proving \eqref{eq:projectivewordtrace} for every $n$.
The denominator two is the actual projector denominator.
In each degree its invariant space has an integral-dimensional
basis of orbit sums of the subset basis; this is not a division
of an unexplained numerical count.

\par\smallskip\noindent\textbf{Centered purity and hard Lefschetz.}\quad
Directly,
\[
 H^{2r}(K_d)=H^{2(r+d)}(X_d)(d)
   =\Q_\ell(-r)^{\oplus\binom{2d}{r+d}},\qquad -d\le r\le d.
\]
The invariant projector acts within each scalar Tate group, so
its image has the same degree-weight relation. This proves the
claimed formula for the cohomology of \eqref{eq:wordcomplex}.
At a point, BBD~5.1.8, pp.~126--127, says that purity of a
complex of weight zero means precisely that its degree-$i$
cohomology is pure of weight $i$ \cite{BBD}.
This is the conclusion here. Its ungraded total has nonzero
weight $-2L$ and $2L$ pieces for $L\ge1$, and is mixed.
Without the centering shifts, the differently twisted dimensions
in \eqref{eq:wordcomplex} need not have the same weight in a
fixed absolute cohomological degree.

Let $\eta_d=\sum_i c_1(\mathcal O_i(1))$.
This is the Chern class of the ample line bundle
$\mathcal O(1,\ldots,1)$, defined over $k$, and $\eta_L$ is
$\sigma$-invariant. Hard Lefschetz applies to the projective
map $X_d\to\Spec k$ and the pure constant intersection complex
\cite[Theorem~5.4.10, p.~144]{BBD}, giving
\begin{equation}\label{eq:wordhardlefschetz}
 \eta_d^{2r}:H^{-2r}(K_d)
          \xrightarrow{\sim} H^{2r}(K_d)(2r),\qquad r\ge0.
\end{equation}
Both this map and its inverse commute with $\sigma$, so it
restricts to an isomorphism on the invariant summand. On
\eqref{eq:wordcomplex}, use the block diagonal operator with
$\eta_L$, $\eta_{L-s}$ and the restriction of $\eta_L$.
All blocks are centered at zero, so the same Lefschetz
isomorphisms hold on their sum.

Below the center, cup product
$H^{2r}\to H^{2r+2}(1)$ is injective: a vector it kills is
also killed by the further power furnishing the Lefschetz
isomorphism to degree $-2r$. Hence $c_r\le c_{r+1}$ for $r<0$,
and \eqref{eq:wordhardlefschetz} gives $c_{-r}=c_r$.
This is the geometric symmetry and unimodality assertion in $v^2$.
In this product example one can also check the Lefschetz action
directly: multiplication by $x$, contraction $x\mapsto1$, and
the degree operator on each $H^*(\PP^1)$ form the standard
$\mathfrak{sl}_2$ representation. Their tensor operators commute
with factor permutations, and the raising operator is cup product
by the displayed invariant ample class.
\end{proof}

\paragraph{Preserving the actual summand is a stronger requirement.}
The trace realization does not extend to every action entry:
$D-2=v^2+v^{-2}$ already has coefficients $(1,0,1)$ in $v^2$.
There is also a basis-independent obstruction. On the actual
rank-two module $N(D-2)$,
\[
 \Tr(P_1P_2)=(D-2)^2=v^4+2+v^{-4}.
\]
This is not an $\mathrm{SL}_2$ character. Any irreducible
constituent supplying weight 4 also supplies weight 2, whereas
the latter coefficient is zero.
Let $e_2$ be the actual commutant projector onto the
multiplicity-four $N(D-2)$ sector of $M$. Then
\[
 \Tr_M(e_2P_1P_2)=4(v^4+2+v^{-4})
\]
is likewise not such a character. The projector belongs to the
regular commutant of \fid{F43}; it is not asserted to belong to
the nonsplit regular right algebra.

Suppose matrix kernels for $P_1,P_2,e_2$ had entries actual
finite-dimensional $\mathrm{SL}_2$ representations, with direct
sum and tensor product as composition, and their character
matrices reproduced these actual operators in some basis.
The sum of the diagonal entries of their composite kernel would
be an actual $\mathrm{SL}_2$ representation. Its character would
equal the last displayed trace, a contradiction. Trace invariance
makes this argument independent of the chosen basis.
Thus that positive kernel framework cannot preserve all three
specified operations simultaneously. This does not exclude the
word-dependent motive just constructed, a model omitting the
projector, virtual kernels, or a category with different operations.

\paragraph{The exact scope of the projective repair.}
The empty word is separately realized by 64 unit objects.
For a nonempty word the total dimension here is
\[
 T_w(1)=20\cdot4^L+20\cdot4^{L-s}.
\]
It varies with the word, consistently with \fid{F42}'s
fixed-common-space obstruction. The present object depends on
$(L,s)$, just as the scalar trace does. We have supplied an actual
projector, polarization and Lefschetz maps, rather than inferred
them from numerical positivity. No action of both nonstandard
generators on these motive objects, compatibility with all word
concatenations, or identification with VIII's designated canonical
basis and structure constants has been proved.

% Separate paragraph for Part III.
\paragraph{\fid{F45} in Grothendieck's language.}
The trace-formula framework is SGA~5, III, 4.7; the data
are smooth projective schemes, rational algebraic
correspondences, a projector, integer Tate twists, a centered
realization and the explicitly chosen ample line bundle
$\mathcal O(1,\ldots,1)$. Its invariant Chern class gives the
Lefschetz isomorphisms on the projector image. Milne's projective
bundle calculation and BBD~5.1.8/5.4.10 supply the exact classical
cohomological framework used above. Cote~16, scans~35, 37 and~39,
is the audited historical anchor for primitive-projector and
Lefschetz language \cite{Archive}; the present polarization and
construction are not attributed to those leaves.
The resulting centered complex is pure of weight zero, with
degree-$i$ cohomology of weight $i$, and its total ungraded
realization has several weights. The equality with the actual
word-trace polynomial is proved for all finite extensions, while
the summand test states exactly which extra compatibility fails
for the indicated positive $\mathrm{SL}_2$-kernel framework.
